\documentclass[12pt,a4paper]{article}

\usepackage{xeCJK}
\usepackage{fontspec}
\usepackage{geometry}
\usepackage{graphicx}
\usepackage{booktabs}
\usepackage{array}
\usepackage{float}
\usepackage[hidelinks]{hyperref}
\usepackage{setspace}
\usepackage{xeCJK}
\usepackage{fontspec}

\setmainfont{Times New Roman}

\geometry{margin=2.5cm}
\setstretch{1.3}

\title{Research Paper Report on Data Center Thermal Management}
\author{Chynna}
\date{\today}

\begin{document}

\maketitle
\tableofcontents
\newpage

\section{ Introduction}

\subsection{Background of Data Center Energy}% 寫資料中心能源消耗、冷卻重要性
With the rapid development of cloud computing and artificial intelligence, the scale and number of data centers have significantly increased, leading to a substantial rise in energy consumption. Among all components, cooling systems account for a large portion of the total energy usage, typically ranging from 30\% to 40\%. To ensure reliable system operation, continuous temperature regulation is required. However, inefficient cooling strategies can result in excessive energy consumption and reduced system efficiency. Therefore, improving energy efficiency, particularly in cooling systems, has become a critical research issue.

\subsection{Importance of Thermal Management}% 寫溫度過高、hotspot、能源浪費
In data centers, servers generate a large amount of heat during operation. Without effective heat dissipation, excessive temperatures and hotspots may occur, which can degrade system performance, shorten equipment lifespan, and increase the risk of failure. On the other hand, overcooling leads to unnecessary energy waste. Therefore, achieving a balance between maintaining safe operating temperatures and minimizing energy consumption is a key challenge in thermal management.

\subsection{Role of AI and Digital Twin}% 寫 AI/ML、Digital Twin 對未來系統的重要性
In recent years, artificial intelligence (AI) and machine learning (ML) techniques have been widely applied to temperature prediction and thermal management in data centers. Data-driven approaches enable more accurate temperature estimation, which helps optimize cooling strategies. In addition, Digital Twin technology provides a virtual representation of the physical system, allowing real-time monitoring and simulation. By integrating AI with Digital Twin, it is possible to improve system efficiency and enable intelligent control.

\subsection{Research Motivation}% 寫你為什麼想研究這個方向
Although existing studies have made progress in temperature prediction, cooling control, and system modeling, most approaches focus on individual components and lack end-to-end integration. For example, some works only address temperature prediction without linking it to control strategies, while others do not consider spatial temperature distribution or real-time implementation. Therefore, this research aims to integrate sensing, prediction, and control mechanisms within a Digital Twin framework to develop a more comprehensive and efficient thermal management system for data centers.

\section{Literature Review}

In recent years, various approaches have been proposed to improve thermal management and energy efficiency in data centers, including temperature prediction, sensor placement, cooling control, and system modeling.

Many studies apply machine learning techniques, such as neural networks and time-series models, to predict temperature dynamics. While these methods effectively capture temporal behavior, they often lack spatial awareness. In addition, sensor placement strategies aim to improve monitoring accuracy, but limited sensors cannot fully represent the overall temperature distribution.

For energy optimization, some works integrate prediction with workload scheduling or cooling control to reduce energy consumption and prevent hotspots. However, these methods are typically designed for specific tasks and lack system-level coordination. 

Recently, Digital Twin technology has been introduced to model and simulate data center systems. Although it enables real-time monitoring and prediction, its integration with control mechanisms is still limited.

Overall, existing studies focus on individual components rather than a unified framework. Therefore, this work aims to explore an integrated approach that combines sensing, prediction, and control for efficient thermal management.




\subsection{Paper 1: Thermal Modeling and Thermal-Aware Energy Saving Methods for Cloud Data Centers}
    
        \subsubsection{Problem}
        
                        The paper addresses the high energy consumption and complexity of thermal management in cloud data centers. Cooling systems consume a significant portion of the total energy, while temperature distribution is difficult to model due to complex airflow and system coupling. Therefore, an efficient thermal modeling and energy optimization strategy is required.

        \subsubsection{Method}

                        This paper is a survey work that categorizes existing approaches into two main parts:
    
         (1) Thermal Modeling
                     
                        \begin{itemize}
                        \item White-box models: Physics-based (e.g., CFD, thermodynamics)
                        \item Black-box models: Data-driven (e.g., ANN, GPR, XGBoost)
                        \item Gray-box models: Combination of physical models and machine learning
                        \end{itemize}
                    
                         
        (2) Energy Saving Methods
      
                        \begin{itemize}
                        \item Thermal-aware IT load scheduling
                        \item Optimization of the cooling system control
                        \item Joint optimization of IT and cooling systems
                        \end{itemize}

        \subsubsection{ Input and Output}

        (1) Input
        
                    The system utilizes various input data, including the following:                    
        \begin{itemize}
                \item Sensor data (temperature, airflow)
                \item IT workload (CPU utilization, power consumption)
                \item Cooling parameters (fan speed, supply temperature)
                \item System state (CRAC, airflow distribution)
        \end{itemize}

        (2) Output
        
                    The output of the system include the following:                   
        \begin{itemize}
            \item Temperature distribution (thermal profile)
            \item Server inlet and outlet temperature
            \item Chip temperature
            \item Energy consumption estimation
            \item Control decisions for scheduling and cooling
        \end{itemize}
        
        \subsubsection{Limitation}

                     The paper has several limitations:      
                \begin{itemize}
                    \item High complexity of thermal modeling due to nonlinear airflow and heat transfer
                \end{itemize}
                \begin{itemize}
                    \item CFD-based methods are accurate but computationally expensive
                    \item Data-driven models depend heavily on data quality
                    \item Weak integration between IT systems and cooling systems
                    \item Liquid cooling technologies are still costly and lack standardization
                \end{itemize}
          
        \subsubsection{Summary}

                This paper provides a comprehensive review of thermal modeling and energy-saving methods in data centers. It highlights the importance of combining physical models with machine learning and emphasizes joint optimization of IT and cooling systems for improving energy efficiency.
\newpage

\subsection{Paper 2: Environmental sensor placement with convolutional }

    \subsubsection{ Problem}
    Environmental sensor placement aims to select optimal locations that maximize information gain and reduce prediction uncertainty, especially in sparse regions like Antarctica. Traditional GP-based methods face limitations in modeling complex non-stationary data and scaling to large datasets.


    \subsubsection{Method}

    This paper proposes a ConvGNP model to learn spatial patterns and uncertainty from both sensor observations and auxiliary data. Sensor locations are selected using uncertainty-based acquisition functions (e.g., JointMI, DeltaVar) to maximize the reduction of prediction uncertainty.

    \subsubsection{ Input and Output}

                \textbf{Input:}
                \begin{itemize}
                    \item Context set 1: Sparse temperature anomaly observations (sensor data)
                    \item Context set 2: Gridded auxiliary data (e.g., elevation, land mask, spatial coordinates)
                    \item Target locations $X^{(t)}$
                \end{itemize}
            
                \textbf{Output:}
                \begin{itemize}
                    \item Predicted temperature anomaly mean $\mu$
                    \item Covariance matrix $K$ representing spatial dependencies
                    \item Suggested sensor placement locations $X^*$
                \end{itemize}

    \subsubsection{Result}

            The ConvGNP significantly outperforms GP-based baselines in terms of RMSE, marginal NLL, and joint NLL. It is able to better capture non-stationary spatial patterns and provide well-calibrated uncertainty estimates. In sensor placement experiments, ConvGNP-based acquisition functions achieve more informative sensor configurations, leading to greater reduction in prediction error compared to traditional methods.
            

    \subsubsection{Limitation}
    
                \begin{itemize}
                    \item The model is trained on simulated reanalysis data (ERA5), which may not fully reflect real-world sensor noise and biases.
                    \item Temporal dependencies are not explicitly modeled, as each time step is treated independently.
                    \item The approach relies on sufficient training data; performance may degrade in extremely sparse observation scenarios.
                    \item The ConvGNP learns non-Bayesian conditioning, which may require large datasets and careful training.
                \end{itemize}
    
    \subsubsection{Summary}

                This paper proposes ConvGNP for sensor placement by modeling spatial uncertainty and selecting locations that maximize information gain, achieving better performance than GP methods.

\newpage

\subsection{Paper 3:
        Intelligent Sensor Placement for Hot Server Detection in Data Centers}
    \subsubsection{Problem}

                Data centers suffer from inefficient cooling due to the lack of accurate thermal monitoring. Existing sensor deployment methods, such as uniform placement, do not consider thermal dynamics, leading to low hot-server detection probability. 
                The key problem is how to optimally place a limited number of sensors to maximize detection performance while controlling false alarms.

    \subsubsection{Method}
    
                The paper proposes a CFD-guided sensor placement approach combined with data fusion and optimization. First, Computational Fluid Dynamics (CFD) is used to model the temperature distribution under various overheating scenarios. Then, a data fusion model is applied to improve detection reliability. Finally, a Lightweight Sensor Placement (LSP) algorithm is developed to efficiently find near-optimal sensor locations that maximize detection probability or minimize the number of sensors.
 
    \subsubsection{Input and Output}

            \textbf{Input:}
            \begin{itemize}
                \item Data center layout (rack positions, CRAC locations)
                \item Airflow and temperature conditions (from CFD simulation)
                \item Sensor noise model and detection threshold
                \item Constraints on false alarm rate and detection probability
            \end{itemize}
            
            \textbf{Output:}
            \begin{itemize}
                \item Optimal sensor placement locations
                \item Detection probability for each monitored location
                \item Minimum number of sensors required to meet performance constraints
            \end{itemize}

    \subsubsection{Result}
    
            Simulation and hardware experiments show that the proposed CFD+LSP method significantly outperforms uniform and random placement strategies. It achieves higher detection probability with fewer sensors. In large-scale data centers, the method can reach over 80\% detection probability with limited sensors, and reduces sensor usage by up to 20\%~50\% compared to baseline approaches.

    
    \subsubsection{Limitation}

            The proposed approach relies heavily on CFD modeling accuracy and requires high computational cost. It assumes static or predefined overheating scenarios, which may not fully reflect dynamic real-world conditions. In addition, practical constraints such as sensor communication, deployment limitations, and hardware variability are not fully considered.

    \subsubsection{Summery}   
            
            This work addresses optimal sensor placement in data centers by leveraging CFD-based thermal modeling and optimization. The proposed CFD+LSP method effectively improves hot-server detection probability while reducing the number of required sensors. Both simulation and hardware results validate its superiority over conventional uniform strategies, although it depends on accurate modeling and has higher computational cost.
\newpage
 
\subsection{Paper 4: Machine Learning-Based Temperature Prediction for Runtime Thermal Management Across System Components}
            
    \subsubsection{Problem}
    
                   \begin{itemize}
                   
                        \item  High temperature causes thermal throttling and performance degradation.
                        \item temperature varies across nodes even under identical workloads.
                        \item Traditional thermal models require hardware-specific information and are not suitable for runtime decision.
                    \end{itemize}
                    
    \subsubsection{Method}

                \begin{itemize}
                   \item Use machine learning models for temperature prediction.
                    \item Models include Reduced Gaussian Process (GP), MLP, and Lasso regression.
                    \item Apply feature selection to reduce complexity.
                    \item Use predicted temperature to guide task placement.
                \end{itemize}
                
    \subsubsection{Input and Output}

                \begin{itemize}
                   \item Input:
                        Application features $A(t)$ (e.g., instructions, cache miss).
                        Physical features $P(t)$ (e.g., temperature, power).
                 \end{itemize}  

                 \begin{itemize}
                   \item Output:
                    Predicted temperature $P(t+1)$.
                    Thermally optimal task placement decision.
                 \end{itemize}
                 
    \subsubsection{Result}

                \begin{itemize}
                    \item Reduced GP reduces the average prediction error from $4.2^\circ C$ to $2.9^\circ C$.
                    \item MLP achieves an average prediction error of about $2.9^\circ C$.
                    \item Lasso regression achieves an average prediction error of about $3.8^\circ C$.
                    \item Task placement success rates are $75\%$ for Reduced GP, $82.5\%$ for MLP, and $74.17\%$ for Lasso.
                    \item The best task placement can reduce temperature by up to $11.9^\circ C$.
                \end{itemize}

    \subsubsection{Limitation}

                \begin{itemize}
                    \item The benchmark diversity is limited.
                    \item The training samples are selected randomly.
                    \item The decoupled model ignores thermal coupling between nodes.
                    \item Large-scale system evaluation is still limited.
                    \item Dynamic scheduling overhead is not fully considered.
                \end{itemize} 
                
    \subsubsection{Summery}
                    This paper proposes a lightweight machine learning-based framework for temperature prediction in HPC systems. By using Reduced Gaussian Process, MLP, and Lasso regression models, the method achieves lower prediction error and reduced computational overhead. The predicted temperature is further applied to thermal-aware task scheduling, enabling better task placement decisions. Among the proposed models, MLP provides the highest success rate in selecting thermally optimal configurations, demonstrating the effectiveness of combining temperature prediction with system-level decision making.

\newpage

\subsection{Paper 5:
        An Intelligent Thermal Management Strategy for a Data Center Prototype Based on Digital Twin Technology}
    \subsubsection{Problem}
    Data centers consume a large amount of energy, where cooling systems account for nearly 50\% of the total energy consumption. Traditional thermal management strategies cannot simultaneously optimize energy efficiency and thermal safety. In addition, limited sensor deployment makes it difficult to reconstruct the full temperature field and detect hotspots in real time.

    \subsubsection{Method}
    This paper proposes a Digital Twin-based thermal management framework. A U-Net model is trained using CFD simulation data to reconstruct the temperature field in real time. A Deep Q-Network (DQN) is applied to optimize the cooling policy by adjusting the CRAC set temperature. The reward function considers PUE reduction, temperature uniformity, and hotspot suppression.

    \subsubsection{Input and Output}
    
                \textbf{Input:}
                \begin{itemize}
                    \item Temperature sensor data (thermocouples)
                    \item Server power (heat load)
                    \item CRAC set temperature
                \end{itemize}
                
                \textbf{Output:}
                \begin{itemize}
                    \item Reconstructed temperature field (U-Net)
                    \item Optimal cooling decision (DQN action)
                \end{itemize}
                
    \subsubsection{Result}
    
    The U-Net achieves a prediction accuracy of MAE = 2.27$^\circ$C and MAPE = 9.21\%. The DQN-based method outperforms fixed-temperature strategies, achieving about 26\% improvement in total reward and 36--40\% improvement in energy efficiency compared to conventional settings. It can also effectively avoid hotspot formation.

    \subsubsection{Limitation}
    The experiment is conducted on a small-scale prototype, which may limit scalability to real data centers. The temperature prediction still contains errors, which may affect decision accuracy. In addition, the DQN training process is time-consuming and may become inefficient in larger and more complex environments.

    \subsubsection{Summery}
    This study integrates Digital Twin, U-Net, and DQN to achieve real-time thermal management in data centers. The proposed method effectively balances energy efficiency and thermal safety, demonstrating the potential of AI-driven control for intelligent data center operation.
\newpage

\subsection{Paper 6: AI-driven-cooling-optimization-in-data-centers-Reinforcement-learning-for-dynamic-workload-placement-and-HVAC-control}
    \subsubsection{Problem}
    This paper focuses on the high cooling energy consumption in data centers. Traditional rule-based workload placement and HVAC control methods cannot adapt well to dynamic workload changes and thermal variations. As a result, they may cause overcooling, thermal hotspots, higher PUE, and increased operational cost.

    \subsubsection{Method}
    The paper proposes a reinforcement learning-based framework to jointly optimize workload placement and HVAC control. The data center is modeled as a dynamic decision-making environment. A multi-agent or hierarchical RL structure is used: one agent manages workload placement, while another agent adjusts HVAC parameters such as cooling setpoints, airflow rates, and fan speeds.

    \subsubsection{Input and Output}
                \textbf{Input:}
                \begin{itemize}
                    \item Temperature data at multiple locations
                    \item Server workload (CPU/GPU utilization)
                    \item HVAC operating parameters (fan speed, airflow, cooling setpoints)
                    \item External environmental conditions (weather, ambient temperature)
                    \item Thermal hotspot indicators or predictions
                \end{itemize}
                
                \textbf{Output:}
                \begin{itemize}
                    \item Workload placement decisions (task migration between servers)
                    \item HVAC control actions (fan speed adjustment, cooling setpoint tuning)
                    \item Airflow and cooling resource allocation
                \end{itemize}

    \subsubsection{Result}
    The proposed RL-based method is evaluated in a simulation environment. Compared with static and heuristic policies, the RL framework reduces cooling energy consumption by about 15\%--22\%. It also improves PUE, reduces thermal hotspot incidents, maintains high SLA compliance, and provides more stable thermal management.

    \subsubsection{Limitation}
    The main limitation is that the evaluation is based on simulation rather than real-world data center deployment. Real data centers have more complex conditions, such as heterogeneous hardware, unpredictable failures, multi-tenant workloads, and changing environmental conditions. In addition, RL training may require high computational cost and careful safety mechanisms before real deployment.

    \subsubsection{Summery}
    This paper demonstrates that reinforcement learning can be used to jointly optimize workload placement and HVAC control in data centers. By considering both computing workload and cooling operation together, the proposed method reduces energy consumption while maintaining thermal safety and SLA performance. However, further real-world validation is still needed to verify its scalability and practical feasibility.
    
\newpage

\subsection{Paper 7:nGRG-RR-2023-01-O-RAN-Towards-6G-v1 3}
    \subsubsection{Problem}
    Current O-RAN architecture is not sufficient to support 6G requirements such as AI-native networking, real-time analytics, energy efficiency, and diverse use cases. There exist gaps between future 6G requirements and existing O-RAN standards.

    \subsubsection{Method}
    A survey-based analysis is conducted by O-RAN members to identify key 6G use cases, architectural gaps, and future directions. The study proposes AI/ML-native architecture, Digital Twin integration, enhanced RAN design, and end-to-end network automation.

    \subsubsection{Input and Output}
                Input:
                Survey responses from O-RAN member companies regarding 6G use cases, architecture, and system requirements.
                
                Output:
                Identified key trends including AI/ML deployment, Digital Twin usage, network automation, security-by-design, and required architectural enhancements for O-RAN towards 6G.

    \subsubsection{Result}
    The study shows that AI/ML in all network layers, real-time analytics, softwarization, and edge-centric architectures are the most critical directions. Most participants agree that O-RAN, 3GPP, and legacy systems will coexist in future 6G networks.

    \subsubsection{Limitation}
    The work is based on a limited number of survey responses and lacks experimental validation, simulation results, and quantitative performance analysis.

    \subsubsection{Summery}
    This report highlights the evolution path of O-RAN towards 6G, emphasizing AI-native architecture, Digital Twin integration, and end-to-end intelligent network management, while identifying key gaps in current systems.
\newpage

\subsection{Paper 8:A Deep Learning-Based Thermal Prediction Approach for Energy Management in Cloud Data Centers}
    \subsubsection{Problem}
                Data centers suffer from high energy consumption and overheating issues due to inefficient VM placement. Existing methods often consider either temperature or power consumption separately, leading to suboptimal thermal management and SLA violations.

    \subsubsection{Method}
                The paper proposes DLTPA, a two-stage framework:
            (1) LSTM is used to predict host temperature and power consumption based on historical data.
            (2) A hybrid GA-SA (Genetic Algorithm + Simulated Annealing) is applied to optimize VM placement, minimizing both temperature and energy consumption.
            
    \subsubsection{Input and Output}
                Input:
            CPU utilization, RAM usage, available RAM, number of CPU cores, used CPU cores, network traffic, number of VMs, power consumption, CPU temperature.
            
            Output:
            Predicted temperature and power consumption, and optimized VM placement strategy.

    \subsubsection{Result}
                LSTM achieves the best prediction performance with high accuracy (R$^2$ up to 0.98).  
            DLTPA reduces peak power consumption by 3.64\%--9.39\%, lowers average temperature by 3.21\%--7.96\%, achieves 0\% SLA violation rate, and maintains good load balancing.

    \subsubsection{Limitation}
                The method is validated mainly in simulation rather than real deployment.  
            It depends heavily on data quality for prediction accuracy.  
            The optimization algorithm (GA-SA) may introduce computational overhead.  
            Only CPU-level thermal and power metrics are considered, without GPU or cooling system modeling.

    \subsubsection{Summery}
    This work integrates deep learning-based prediction with heuristic optimization for thermal-aware VM placement. It effectively reduces temperature and energy consumption while maintaining SLA, demonstrating the advantage of combining prediction and decision-making in data center management.
\newpage

\subsection{Paper 9:Energy-Efficient Task Distribution Using Neural Network Temperature Prediction in a Data Center}
    \subsubsection{Problem}
    Data centers suffer from high energy consumption, especially due to cooling systems. Traditional load balancing methods only consider CPU utilization or current temperature, ignoring future temperature rise. This may lead to thermal hotspots and inefficient energy usage.


    \subsubsection{Method}
    This paper proposes a neural network (NN)-based temperature prediction model to estimate future CPU temperature (1-second ahead). The predicted temperature is then used to guide task distribution, aiming to reduce the maximum CPU temperature across servers.
    
    \subsubsection{Input and Output}
                    Input:
                \begin{itemize}
                    \item Current CPU temperature
                    \item CPU usage
                    \item Local fan speed
                    \item Number of CPU cores requested by tasks
                \end{itemize}
                
                Output:
                \begin{itemize}
                    \item Predicted future CPU temperature
                \end{itemize}
                
    \subsubsection{Result}
    The NN model achieves accurate temperature prediction with low MSE. Simulation results show that the proposed method can reduce the maximum CPU temperature (up to $\sim 0.45^\circ C$) compared to non-predictive methods, especially when multiple allocation trials are considered.

    \subsubsection{Limitation}
                \begin{itemize}
                \item Improvement in temperature reduction is limited
                \item High computational cost (not suitable for real-time control)
                \item Does not consider airflow or inlet temperature variations
                \item Evaluated on a small number of servers
                \item Task distribution evaluated only in simulation
                \end{itemize}
                
    \subsubsection{Summery}
    This work integrates temperature prediction with task distribution, demonstrating that considering future thermal behavior can improve load balancing. However, practical deployment requires faster computation and more comprehensive modeling.

\newpage
\subsection{Paper 10:Integrating Experimental Numerical and Machine Learning Models for Real-Time Efficient Data Center Cooling Control}
    \subsubsection{Problem}
    Traditional air-cooled data centers mainly use reactive cooling control, which adjusts the cooling resources only after temperature increases are detected. This leads to slow response time, overcooling, energy waste, and difficulty handling dynamic workload variations. In addition, CFD-based thermal modeling is computationally expensive and unsuitable for real-time cooling control.

    \subsubsection{Method}
    The paper proposes a Dynamic Cooling System Control (DCSC) mechanism integrating experimental data, machine learning models, and thermal models for predictive cooling control. Real-time data are collected using an aggregator system. A feed-forward neural network is trained to estimate CPU temperatures based on server utilization, CRAH settings, and pressure sensor data. A power predictor forecasts future server power consumption, while a compact thermal model serves as a backup model for robustness.

    \subsubsection{Input and Output}
                \begin{itemize}
                \item Input:
                \begin{itemize}
                    \item CPU A utilization
                    \item CPU B utilization
                    \item CRAH 1 and CRAH 2 VFD speeds
                    \item CRAH supply air temperatures
                    \item Cold aisle pressure sensor data
                    \item Plenum pressure sensor data
                \end{itemize}
            
                \item Output:
                \begin{itemize}
                    \item Predicted CPU A temperature
                    \item Predicted CPU B temperature
                    \item Cooling control information for CRAH units
                \end{itemize}
            \end{itemize}

    \subsubsection{Result}
    The proposed feed-forward neural network achieved high prediction accuracy for CPU temperature estimation. The testing results showed MAE values of approximately $0.85^\circ C$ for CPU A and $0.81^\circ C$ for CPU B, with $R^2$ values close to 0.985. The average prediction latency was only $75\,\mu s$ per sample, demonstrating the feasibility of real-time cooling prediction. Experimental analysis also revealed airflow imbalance and hotspot formation inside the aisle, highlighting the importance of zone-based thermal management.

    \subsubsection{Limitation}
                \begin{itemize}
                \item The proposed system was only validated on a limited zone of the data center.
                \item The complete closed-loop cooling control system was not fully implemented.
                \item The experiments mainly focused on controlled operating conditions.
                \item The scalability to different data center layouts was not fully verified.
                \item Sudden workload spikes and complex real-world disturbances were not extensively evaluated.
                \end{itemize}
% Paper 3 ~ Paper 10 照這個格式複製
\newpage

\section{Category-based Analysis}

    \subsection{Sensing Methods}
    Sensing methods are used to collect environmental and operational data in data centers, such as temperature, airflow, and server utilization. Many studies focus on optimizing sensor placement to improve monitoring accuracy under limited sensing resources. These methods can obtain real measurements with relatively low cost. However, sensor data are typically localized and cannot represent the global temperature distribution. Therefore, sensing alone is insufficient for comprehensive thermal management.

    \subsection{Temperature Prediction Methods}
    Temperature prediction methods usually employ machine learning models, such as neural networks and LSTM, to estimate future temperature based on historical data and system workload. These approaches are effective in capturing temporal dynamics and can provide early warnings of temperature increase. However, most methods focus on point-wise prediction and lack spatial awareness, making them insufficient for detecting hotspots across the entire data center.

    \subsection{Temperature Reconstruction Methods}
    Temperature reconstruction methods aim to recover the full temperature field from limited sensor data. Deep learning models, such as convolutional neural networks and U-Net, are commonly used to capture spatial relationships. These approaches can provide a global view of temperature distribution and improve hotspot detection. However, they usually require a large amount of training data and may suffer from limited generalization in real-world environments.

    \subsection{Control and Optimization Methods}
    Control and optimization methods utilize sensing or prediction results to adjust system operations, such as workload scheduling and cooling control. Common techniques include reinforcement learning and genetic algorithms. These methods can effectively reduce energy consumption and prevent overheating. However, most of them are designed for specific objectives and lack coordination across different system components.

    \subsection{System-level Modeling and Digital Twin}
    System-level modeling methods aim to construct a comprehensive representation of the data center, such as Digital Twin models. These approaches enable simulation, prediction, and real-time decision-making by integrating multiple system components. While Digital Twin has the potential to unify sensing, prediction, and control, its practical implementation still faces challenges in accuracy and real-time performance.

\section{Cross-paper Comparison}

    \subsection{Overall Comparison}
            \begin{table}[H]
                    \centering
                    \small
                    \caption{Comparison of Reviewed Papers}
                        \begin{tabular}{p{2cm} p{3cm} p{3cm} p{3cm} p{3cm}}
                        \toprule
                        Paper & Method & Input & Output & Limitation \\
                        \midrule
                        Paper 1 & Review of thermal modeling and energy-saving methods & Existing literature, system models & Categorized methods and research gaps & No experimental validation, lacks unified framework \\
                        
                        Paper 2 & ConvGNP + active learning for sensor placement & Sensor data, spatial features (elevation, mask) & Optimal sensor locations, probabilistic prediction & Uses simulated data, no temporal modeling \\
                        
                        Paper 3 & CFD-based modeling + optimization for sensor placement & Server layout, airflow, thermal data & Optimal sensor placement, improved hotspot detection & High computational cost, depends on CFD accuracy \\
                        
                        Paper 4 & ML-based prediction (GP, NN, Lasso) for scheduling & Application features, physical features & Temperature prediction, task scheduling decision & Requires training data, lacks spatial modeling \\
                        
                        Paper 5 & Digital Twin + U-Net + DQN & Sensor data, server power, AC temperature & Temperature field reconstruction, optimal cooling policy & Small-scale prototype, relies on CFD data \\
                        
                        Paper 6 & Reinforcement learning (PPO, DQN) for workload placement and HVAC control & Temperature, server load, HVAC parameters, external conditions & Optimal workload distribution and cooling policy (reduced PUE, improved thermal stability) & Simulation-based validation, high training cost, limited real-world deployment \\
                        
                        Paper 7 & Survey and architecture analysis for AI-native 6G O-RAN systems & Survey data, use cases, network requirements & Future 6G architecture, AI/ML integration, system design directions & No implementation, lacks quantitative evaluation, high-level concepts only \\
                        
                        Paper 8 & LSTM-based thermal prediction + GA-SA optimization for VM placement & CPU utilization, memory, network I/O, number of VMs, historical temperature and power data & Predicted temperature and power consumption, optimized VM placement & Simulation-based evaluation, no real-time deployment, high model and optimization complexity\\
                        
                        Paper 9 & Neural network-based temperature prediction for task distribution & Current CPU temperature, fan speed, CPU usage, requested CPU cores & Predicted future CPU temperature and task allocation decisions & High computational cost, limited temperature reduction, short-term prediction only (1s ahead) \\
                        
                        Paper 10 &  &  &  &  \\
                        \bottomrule
                        \end{tabular}
                \end{table}

    \subsection{Method Comparison}
                Table~\ref{tab:method_comparison} summarizes the main methodological differences among the reviewed papers.

                \begin{table}[H]
                    \centering
                    \small
                    \caption{Comparison of Method Types}
                    \label{tab:method_comparison}
                    \begin{tabular}{p{3cm} p{4cm} p{4cm} p{4cm}}
                    \toprule
                         Method Type & Related Papers & Strength & Limitation \\
                    \midrule
                    Survey / Review & Paper 1, Paper 7 & Provides background, taxonomy, and future trends & No implementation or quantitative evaluation \\
                    
                    Physics-based Modeling & Paper 3 & High accuracy for airflow and thermal behavior & High computational cost; not suitable for real-time control \\
                    
                    Data-driven Prediction & Paper 2, Paper 4, Paper 8, Paper 9 & Fast prediction and useful for dynamic decision-making & Depends on data quality and lacks physical interpretability \\
                    
                    Optimization / Control & Paper 6, Paper 8, Paper 9 & Can reduce temperature, power, or improve workload placement & Usually focuses on specific objectives only \\
                    
                    System-level / Hybrid Framework & Paper 5, Paper 10 & Integrates sensing, prediction, and control & More complex and difficult to deploy in real data centers \\
                    \bottomrule
                    \end{tabular}
                \end{table}
                
    \subsection{Static vs Dynamic Analysis}

            \begin{table}[H]
            \centering
            
            \caption{Comparison between Static and Dynamic Methods}
            \label{tab:static_dynamic}
            \begin{tabular}{p{3cm} p{5cm} p{5cm}}
            \toprule
             & Static Methods & Dynamic Methods \\
            \midrule
            Definition & One-time analysis without time evolution & Time-dependent modeling and real-time updates \\
            
            Representative Papers & Paper 1, Paper 2 (partial), Paper 3 & Paper 4, Paper 6, Paper 8, Paper 9, Paper 10 \\
            
            Strength & High accuracy (especially physics-based models) & Real-time adaptability and responsiveness \\
            
            Limitation & Cannot adapt to workload changes, not suitable for real-time control & Prediction error accumulation, depends on data quality \\
            
            Use Case & System design and offline analysis & Runtime thermal management and control \\
            \bottomrule
            \end{tabular}
            \end{table}
            
    \subsection{Real-world Deployment Analysis}

                    \begin{table}[H]
                    \centering
                    \caption{Real-world Deployment Analysis}
                    \label{tab:deployment_analysis}
                    \begin{tabular}{p{3cm} p{4cm} p{4cm} p{4cm}}
                    \toprule
                    Method Type & Related Papers & Deployment Advantage & Deployment Challenge \\
                    \midrule
                    Review & Paper 1, Paper 7 & Provides research direction and system requirements & No direct implementation or validation \\
                    
                    Sensor Placement & Paper 2, Paper 3 & Improves monitoring accuracy and supports hotspot detection & Sensor cost, placement constraints, and environmental changes \\
                    
                    Prediction Models & Paper 4, Paper 8, Paper 9 & Fast prediction for runtime thermal management & Requires continuous data collection and model updating \\
                    
                    RL-based Control & Paper 6 & Can adapt cooling and workload decisions dynamically & Needs safe training, stable simulation, and deployment safeguards \\
                    
                    Digital Twin / Hybrid Framework & Paper 5, Paper 10 & Supports real-time visualization, prediction, and control integration & High system complexity and difficult scalability \\
                    \bottomrule
                    \end{tabular}
                    \end{table}
                    
                    \begin{itemize}
                        \item Prediction-based methods are relatively easier to deploy because they can be integrated with existing monitoring systems.
                        
                        \item Sensor placement methods are practical but may require hardware modification and additional installation cost.
                        
                        \item RL-based control methods have strong potential, but direct deployment in real data centers is risky without safety mechanisms.
                        
                        \item Digital Twin and hybrid frameworks provide the most complete solution, but they require high integration effort, real-time data pipelines, and reliable system models.
                        
                        \item Overall, real-world deployment requires not only prediction accuracy, but also safety, scalability, and compatibility with existing data center infrastructure.
                    \end{itemize}
                    
\section{Research Gap}

    \subsection{Lack of End-to-End Integration}
                Most existing studies focus on individual components of data center thermal management, such as sensing (Paper 2, 3), prediction (Paper 4, 8, 9), or control (Paper 6). Although system-level approaches (Paper 5 and Paper 10) provide partial integration, these modules are often developed independently.
        
                As a result, there is still a lack of a unified framework that connects sensing, prediction, decision-making, and control. This limitation makes it difficult to achieve efficient and real-time thermal management in practical data centers.
        
                Therefore, developing an end-to-end integrated system remains a key research gap.
    
    \subsection{Lack of Spatial Temperature Modeling}
                \begin{itemize}
                \item Many studies only perform point-wise temperature prediction (Paper 4, Paper 8, Paper 9).
                
                \item These methods cannot capture the overall spatial temperature distribution of data centers.
                
                \item CFD-based methods (Paper 3) provide spatial information but require high computational cost.
                
                \item U-Net-based approaches (Paper 5) improve spatial modeling, but real-time deployment remains challenging.
                
                \item As a result, insufficient spatial modeling limits hotspot detection and precise thermal control.
                \end{itemize}
                
    \subsection{Limited Digital Twin-based Control}
                \begin{itemize}
                \item Digital Twin frameworks (Paper 5 and Paper 7) provide virtual system modeling and real-time monitoring.
                
                \item However, most studies focus on visualization and prediction rather than closed-loop thermal control.
                
                \item Existing methods have limited integration between Digital Twin, AI prediction, and control strategies.
                
                \item In addition, real-world deployment remains challenging due to system complexity and computational overhead.
                
                \item Therefore, intelligent control based on Digital Twin remains an open research problem.
            \end{itemize}
            
    \subsection{Real-time Implementation Challenges}
                \begin{itemize}
                \item Many thermal management methods are validated only in simulation environments.
                
                \item Physics-based approaches, such as CFD modeling (Paper 3), require high computational cost and are difficult to apply in real time.
                
                \item Machine learning methods (Paper 4, Paper 8, and Paper 9) provide faster prediction, but their accuracy depends heavily on training data.
                
                \item Digital Twin and hybrid frameworks (Paper 5 and Paper 10) improve system integration, but real-time deployment still faces challenges in scalability and latency.
                
                \item Therefore, achieving accurate and efficient real-time thermal management remains a major challenge in practical data centers.
                \end{itemize}


\section{Conclusion}

    After reviewing the ten papers, it can be observed that data center cooling systems are evolving toward predictive and intelligent thermal management using machine learning, Digital Twin, and real-time monitoring techniques.However, most studies only focus on individual tasks such as temperature prediction or cooling control. Therefore, this research aims to combine Digital Twin and neural networks to develop a real-time cooling optimization framework for improving energy efficiency and reducing hotspots.


\end{document}